\chapter{Conclusioni e sviluppi futuri}
\label{sec:conclusion}

\section{Sintesi dei risultati}
\label{sec:concl-summary}

Questa tesi si è chiesta se un modello linguistico compatto, specializzato in un'unica decisione e impiegato come
decision model, possa verificare frase per frase, su un portatile, le risposte di un sistema RAG, e da dove provenga
l'eventuale vantaggio. GroundOne, un modello Qwen2.5 sottoposto a fine-tuning che legge i verdetti di un'intera risposta
in un solo forward pass, fornisce una risposta positiva, sebbene con alcune riserve.

La specializzazione funziona là dove era mirata. Sul test set di RAGTruth GroundOne raggiunge una balanced accuracy a tre
classi di 0.769 con 0.5B parametri e di 0.832 con 3B. Supera di molto gli stessi modelli non addestrati, ordina le frasi
supportate meglio degli encoder specializzati ed è nettamente superiore a Rizzo-Flow, un decision model generalista otto
volte più grande. Su LLM-AggreFact il 3B è alla pari con il miglior encoder di piccole dimensioni. Su VitaminC, invece,
il guadagno scompare: uno specialista non è uno strumento generale di verifica dei fatti. d1-3B, un decision model
generalista della stessa dimensione addestrato a decidere su un'ampia miscela di task, conferma il compromesso: GroundOne
lo supera nettamente su RAGTruth, lo eguaglia su LLM-AggreFact e resta molto indietro su VitaminC.

Il meccanismo del decision model ripaga su due fronti. Il primo è la velocità. Leggere tutti i verdetti in un solo forward pass
rende la verifica delle risposte lunghe circa due volte più rapida di quella di un giudice generativo sintetico e sei
volte più rapida di quella di un giudice verboso della stessa dimensione, con un output sempre valido. Il secondo è
l'accuratezza: restringere la decisione alle risposte ammissibili è di per sé più accurato che lasciare allo stesso
modello la generazione del verdetto. Inoltre, anche quando si chiede loro lo stesso output, i giudici generativi non
riescono a fornire probabilità di qualità paragonabile. Le loro probabilità separano a malapena le decisioni corrette da
quelle errate, e quelle espresse a parole sono quasi sempre del tipo tutto-o-niente.

Le probabilità di GroundOne si possono calibrare con una temperatura e un bias stimati sul dev set. Una soglia scelta
sugli stessi dati le converte poi in etichette coerenti sia con le probabilità sia con la decisione valutata.
L'Expected Calibration Error, da solo, non distingue però le probabilità utili da quelle inutili: la capacità
discriminativa va misurata a parte.

L'intero sistema gira su un portatile da 16\,GB. Una risposta lunga viene verificata in 0.41 secondi e 5.5 joule con il
0.5B, e in 2.57 secondi e 33 joule con il 3B. Rispetto a un giudice generativo il risparmio di energia è reale, ma
inferiore al risparmio di tempo, perché un singolo prefill lungo assorbe più potenza del decoding.

I limiti principali sono due. Il primo è la ristrettezza dei dati di addestramento, a causa della quale le
contraddizioni al di fuori del data-to-text e le allucinazioni nascoste in frasi copiate vengono apprese solo in parte.
Il secondo è una ricetta di addestramento non del tutto robusta rispetto alla scelta del seed.

\section{Sviluppi futuri}
\label{sec:concl-future}

Da questi limiti discendono direttamente diverse linee di lavoro. La ricetta di addestramento andrebbe resa robusta
abilitando la bias correction e il gradient clipping, registrando le norme dei gradienti e ripetendo i run con più seed,
anche per il 3B. Servirebbero poi dati di addestramento con contraddizioni più varie, con supporto parafrasato e con
modifiche minime come quelle di VitaminC: dovrebbero intervenire sia sulla bassa recall al di fuori del data-to-text sia
sul punto cieco delle frasi copiate. Infine, addestrare il readout sulla lettera che il modello scriverebbe
spontaneamente dopo il marker potrebbe allineare readout e generazione, con effetti sull'accuratezza ancora da misurare.

Sul versante dei sistemi, il modello di costo del Capitolo~\ref{sec:local} andrebbe verificato su altro hardware. Un
porting del readout su llama.cpp permetterebbe di misurare GroundOne su una scheda grafica dedicata, su un DGX Spark o
su una macchina AMD Strix Halo, e su uno smartphone, le cui unità di elaborazione neurale accelerano proprio il prefill
su cui si regge un decision model. Per il pilastro della località resta inoltre da misurare la memoria del retriever accanto a
generatore e verificatore, già misurati insieme in questa tesi. Anche i confronti di latenza con le altre famiglie di verificatori richiedono condizioni eque: gli encoder
vanno eseguiti in batch, mentre GroundOne e i decision model generalisti vanno interrogati attraverso la stessa interfaccia
HTTP.

La verifica potrebbe inoltre procedere durante la generazione della risposta, segnalando ogni frase non appena è
completa: la soluzione basata sui marker rende naturale questo uso. Infine, un assistente locale costruito
sull'interfaccia del servizio permetterebbe agli utenti di scegliere la regola e la soglia più adatte ai propri
documenti. È proprio in questo uso che le probabilità calibrate mostrano il loro valore.
